% !TEX root = ../main.tex
\chapter*{How to read this manual}
\addcontentsline{toc}{chapter}{How to read this manual}
\markboth{How to read this manual}{}

This manual describes \emph{Underground Brains}, a web application for looking at and
analysing long electrical recordings from fungi. It is written for two kinds of reader,
and laid out so that each can skip what the other needs.

\begin{plainterms}[Two readers, one book]
Think of a museum. Every exhibit has a placard for visitors and a catalogue entry for
specialists, and each points to the other. Part~I is the placards: how to use the site, and
why, in everyday words. Part~II and the appendices are the catalogue: what the code does,
exactly enough to check it or to build on it.
\end{plainterms}

\section*{If you want to use the site}
Read \cref{part:using}. Start with \cref{ch:why}, then \cref{ch:first-session}, which takes
one recording from start to finish in about fifteen minutes, then \cref{ch:core-ideas},
which explains the handful of words the rest of the manual leans on.

The chapters after that follow the path the app draws across the top of every page:
\begin{center}\sffamily
\textbf{Explore}: find a span \enskip\ensuremath{\to}\enskip
\textbf{Analyse}: run a chain over it \enskip\ensuremath{\to}\enskip
\textbf{Review}: judge what it found
\end{center}
with \textbf{Library} and \textbf{Discovery} feeding Analyse, and \textbf{Models} reading
what it produced. \Cref{ch:workflows} shows complete pieces of research done with the app,
\cref{ch:faq} answers common questions, and \cref{ch:glossary} defines every term.

\section*{If you want to know exactly what it does}
Every analysis step in the app is a \emph{block}. \Cref{ch:blocks} gives each of the
\ubnumblocks{} blocks a short card, and every card links to the block's full entry in
\cref{app:blocks}: the algorithm, its assumptions, every setting, how it fails, how it was
checked, and where its code and tests live. \Cref{part:how} explains how the app is built
and how to extend it; \cref{app:nulls,app:data} cover the statistics the app uses to tell a
pattern from chance, and how recordings are stored.

\section*{Signs used in this manual}
\begin{plainterms}
A blue box explains a technical idea in everyday words.
\end{plainterms}
\begin{watchout}
An amber box warns about something that costs time or gives a wrong answer.
\end{watchout}
\begin{underthehood}
A grey box points into the code. Readers who only use the site can skip these.
\end{underthehood}

Each page of the app is described in its own section, marked \pagestatus{live} when it works
on your own recordings, \pagestatus{preview} when it still shows example data, or
\pagestatus{unchecked} until someone has checked. Each block card carries the coloured tag of
the tab it sits under in the app's Insert-a-stage window, such as \catbadge{detect} or
\catbadge{cluster}.

\section*{Where this manual comes from}
The manual lives in the code repository, in \code{docs/manual/}, next to the code it
describes. The list of blocks and the settings table in every block entry are generated from
the code by \code{scripts/manual_tables.py}, so they cannot drift from it; everything else is
written by hand. This copy documents the code at commit \texttt{\ubversion}.

While the manual is a draft, an unfinished passage is marked
\towritesample{like this}, and a claim nobody has yet checked against the code or a source
is marked \unverifiedsample{like this}. The title page counts both: \ubwritingstatus.
